%! TEX root = part6.tex
% the main TeX file which is intended to compile, :VimtexReload after adjustment
% :h vimtex-tex-root
\documentclass[../main.tex]{subfiles}
% \includeonlyframes{current}
% \setbeameroption{hide notes} % Only slides
% \setbeameroption{show only notes} % Only notes
\setbeameroption{show notes on second screen=right} % Both

\begin{document}

\begin{frame}[mytitle=standard,c,mycolor=digiPH_white,light]
	\frametitle{Tugas Kelas}

	\normalsize	Buatlah Gambar Kerja \huge Rencana Balok \normalsize

\end{frame}


\begin{frame}[mytitle=standard,t,mycolor=digiPH_uniblue,light]\frametitle{Struktur Beton}
	\begin{itemize}
		\item Beton adalah campuran semen, agregat, air, dan bahan tambahan
		\item Komposisi umum:
		      \begin{itemize}
			      \item ±15\% semen
			      \item ±8\% air
			      \item ±3\% udara
			      \item Sisanya agregat
		      \end{itemize}
		\item Kuat tekan tinggi
		\item Mencapai kekuatan rencana pada umur 28 hari
	\end{itemize}
\end{frame}

\begin{frame}[mytitle=standard,t,mybg=bglight2,light]\frametitle{Karakteristik Beton}
	\begin{itemize}
		\item Kuat terhadap tekan
		\item Lemah terhadap tarik → perlu tulangan baja
		\item Tahan terhadap lingkungan (air, api, korosi)
		\item Mudah dibentuk (bekisting)
	\end{itemize}\end{frame}

\begin{frame}[mytitle=standard,t,mybg=bglight3,light]\frametitle{Bahan Tambahan Beton}
	\begin{itemize}
		\item Fly ash
		\item Slag
		\item Silica fume
		\item Dust collector
		\item Fungsi:
		      \begin{itemize}
			      \item Meningkatkan workability
			      \item Menambah durabilitas
			      \item Mengatur waktu pengerasan
		      \end{itemize}
	\end{itemize}\end{frame}

\begin{frame}[mytitle=standard,t,mybg=bglight3,light]\frametitle{Jenis Beton}
	\begin{itemize}
		\item Beton bertulang
		      \begin{itemize}
			      \item Menggunakan tulangan baja
			      \item Menahan tarik dan lentur
		      \end{itemize}
		\item Beton pratekan
		      \begin{itemize}
			      \item Diberi tegangan awal
			      \item Mengurangi retak
		      \end{itemize}
		\item Beton pracetak
		      \begin{itemize}
			      \item Dicetak di pabrik
			      \item Kualitas lebih terkontrol
		      \end{itemize}
	\end{itemize}\end{frame}

\begin{frame}[mytitle=standard,t,mybg=bglight2,light]\frametitle{Keunggulan Beton}
	\begin{itemize}
		\item Workability tinggi
		\item Kuat tekan besar
		\item Tahan panas dan lingkungan
		\item Biaya relatif murah
	\end{itemize}\end{frame}

\begin{frame}[mytitle=standard,t,mycolor=digiPH_cream,light]\frametitle{Material Penyusun Beton}
	\begin{itemize}
		\item Semen → bahan pengikat
		\item Agregat kasar → kekuatan utama
		\item Agregat halus → mengisi celah
		\item Air → reaksi hidrasi
		\item Admixture → meningkatkan performa
	\end{itemize}\end{frame}

\begin{frame}[mytitle=standard,t,mybg=bglight3,light]\frametitle{Jenis Semen}

	\ltxt{digiPH_white}{digiPH_black}{\begin{itemize}
			\item OPC (Ordinary Portland Cement)
			\item PPC (Portland Pozzolan Cement)
			\item Sifat penting:
			      \begin{itemize}
				      \item Kehalusan (fineness)
				      \item Mempengaruhi kecepatan hidrasi
			      \end{itemize}
		\end{itemize}}
	\rimg{jenis_semen.jpeg}

\end{frame}

\begin{frame}[mytitle=standard,t,mycolor=digiPH_gray,light]\frametitle{Agregat dan Air}
	\begin{itemize}
		\item Agregat kasar:
		      \begin{itemize}
			      \item Kerikil, batu pecah
		      \end{itemize}
		\item Agregat halus:
		      \begin{itemize}
			      \item Pasir
		      \end{itemize}
		\item Air:
		      \begin{itemize}
			      \item Terlalu banyak → beton lemah
			      \item Terlalu sedikit → sulit dikerjakan
		      \end{itemize}
	\end{itemize}\end{frame}

\begin{frame}[mytitle=standard,t,mycolor=digiPH_purple,dark]\frametitle{Admixture Beton}
	\begin{itemize}
		\item Superplasticizer → meningkatkan workability
		\item Retarder → memperlambat pengerasan
		\item Accelerator → mempercepat pengerasan
		\item Air entraining → meningkatkan durabilitas
		\item Waterproofing → mengurangi permeabilitas
	\end{itemize}\end{frame}

\begin{frame}[mytitle=standard,t,mycolor=digiPH_ubcblue,dark]\frametitle{Prinsip Desain Struktur Beton}
	\begin{itemize}
		\item Kekuatan
		\item Stabilitas
		\item Kekakuan
		\item Durabilitas
		\item Keamanan
	\end{itemize}\end{frame}

\begin{frame}[mytitle=standard,t,mycolor=digiPH_uniblue,light]\frametitle{Desain Balok}
	\begin{itemize}
		\item Menahan beban lentur
		\item Langkah desain:
		      \begin{itemize}
			      \item Hitung beban dan momen
			      \item Tentukan dimensi
			      \item Penempatan tulangan
		      \end{itemize}
	\end{itemize}\end{frame}

\begin{frame}[mytitle=standard,t,mycolor=digiPH_red,dark]\frametitle{Desain Kolom}
	\begin{itemize}
		\item Menahan beban tekan
		\item Langkah:
		      \begin{itemize}
			      \item Hitung beban vertikal
			      \item Tentukan dimensi
			      \item Penulangan
		      \end{itemize}
	\end{itemize}\end{frame}

\begin{frame}[mytitle=standard,t,mycolor=digiPH_leaf,dark]\frametitle{Desain Plat/Slab}
	\begin{itemize}
		\item Elemen datar penahan beban
		\item Perlu:
		      \begin{itemize}
			      \item Analisis lentur
			      \item Distribusi beban
			      \item Penempatan tulangan
		      \end{itemize}
	\end{itemize}\end{frame}

\begin{frame}[mytitle=standard,c,mycolor=digiPH_uniblue,light]\frametitle{Beton Pratekan}
	\begin{itemize}
		\item Baja ditarik sebelum pengecoran
		\item Setelah beton mengeras → gaya tekan
		\item Keunggulan:
		      \begin{itemize}
			      \item Kapasitas lentur tinggi
			      \item Retak lebih kecil
		      \end{itemize}
	\end{itemize}
\end{frame}

\begin{frame}[mytitle=standard,c,mycolor=digiPH_cream,light]\frametitle{Jenis Beban pada Beton}
	\begin{itemize}
		\item Beban statis (tetap)
		\item Beban dinamis (berubah)
		\item Beban hidup (aktivitas)
		\item Beban mati (berat struktur)
	\end{itemize}\end{frame}

\begin{frame}[mytitle=standard,t,mybg=bglight1,light]\frametitle{Analisis Struktur Beton}
	\limg{analisis_beton.jpeg}

	\rtxt{digiPH_gray}{digiPH_darkblue}{\begin{itemize}
			\item Digunakan untuk:
			      \begin{itemize}
				      \item Menentukan respon struktur
				      \item Menghitung tegangan dan deformasi
			      \end{itemize}
			\item Metode analisis diperlukan untuk keamanan struktur
		\end{itemize}}



\end{frame}

\begin{frame}[mytitle=standard,t,mycolor=digiPH_purple,dark]
	\frametitle{Kerusakan Beton}
	\rtxt{digiPH_white}{digiPH_black}{\begin{itemize}
			\item Faktor:
			      \begin{itemize}
				      \item Lingkungan
				      \item Mekanik
				      \item Kesalahan konstruksi
				      \item Usia struktur
			      \end{itemize}
			\item Bentuk:
			      \begin{itemize}
				      \item Retak
				      \item Keropos
				      \item Spalling
			      \end{itemize}
		\end{itemize}}
	\limg{rusak_beton.jpeg}
\end{frame}

\begin{frame}[mytitle=standard,t,mybg=bglight3,light]\frametitle{Perbaikan Beton}
	\begin{itemize}
		\item Rekonstruksi parsial
		\item Penguatan tulangan
		\item FRP
		\item Grouting dan epoxy injection
		\item Shotcrete dan patching
	\end{itemize}\end{frame}

\begin{frame}[mytitle=center,c,mycolor=digiPH_white,light]\frametitle{Inovasi Beton}
	\begin{itemize}
		\item Smart concrete (sensor)
		\item Beton permeabel
		\item Beton ramah lingkungan
		\item Teknologi 3D printing
	\end{itemize}\end{frame}



\end{document}
